\documentclass[12pt,a4paper]{article}

% Encoding and fonts
\usepackage[utf8]{inputenc}
\usepackage[T1]{fontenc}
\usepackage{lmodern}

% Page layout
\usepackage[top=1in, bottom=1in, left=1in, right=1in]{geometry}

% Typography
\usepackage{microtype}
\usepackage{parskip}

% Language
\usepackage[english]{babel}

% Headers and footers
\usepackage{fancyhdr}
\pagestyle{fancy}
\fancyhf{}
\fancyhead[L]{\leftmark}
\fancyhead[R]{\thepage}
\fancyfoot[C]{\thepage}
\renewcommand{\headrulewidth}{0.4pt}
\renewcommand{\footrulewidth}{0pt}

% Section styling
\usepackage{titlesec}
\titleformat{\section}{\Large\bfseries}{\thesection}{1em}{}
\titleformat{\subsection}{\large\bfseries}{\thesubsection}{1em}{}
\titleformat{\subsubsection}{\normalsize\bfseries}{\thesubsubsection}{1em}{}
\titlespacing*{\section}{0pt}{2.5ex plus 1ex minus .2ex}{1.5ex plus .2ex}
\titlespacing*{\subsection}{0pt}{2ex plus 1ex minus .2ex}{1ex plus .2ex}
\titlespacing*{\subsubsection}{0pt}{1.5ex plus 1ex minus .2ex}{0.8ex plus .2ex}

% List formatting
\usepackage{enumitem}
\setlist[itemize]{topsep=0.5ex, itemsep=0.25ex, parsep=0pt}
\setlist[enumerate]{topsep=0.5ex, itemsep=0.25ex, parsep=0pt}

% Essential packages
\usepackage{graphicx}
\usepackage[table]{xcolor}
\usepackage{booktabs}
\usepackage{array}
\usepackage{tabularx}
\usepackage{longtable}
\usepackage{amsmath}
\usepackage{amssymb}
\usepackage[normalem]{ulem}
\rowcolors{2}{gray!10}{white}
\renewcommand{\arraystretch}{1.3}

% Cross-references
\usepackage{hyperref}
\usepackage{cleveref}

% Hyperref setup
\hypersetup{
    colorlinks=true,
    linkcolor=blue,
    filecolor=magenta,
    urlcolor=cyan,
}

% Code listings
\usepackage{listings}
\definecolor{codebg}{RGB}{248,248,248}
\definecolor{codeframe}{RGB}{200,200,200}
\definecolor{codekeyword}{RGB}{0,0,180}
\definecolor{codecomment}{RGB}{0,128,0}
\definecolor{codestring}{RGB}{163,21,21}
\lstset{
    basicstyle=\ttfamily\small,
    breaklines=true,
    breakatwhitespace=false,
    frame=single,
    framerule=0.5pt,
    rulecolor=\color{codeframe},
    backgroundcolor=\color{codebg},
    keepspaces=true,
    showstringspaces=false,
    tabsize=4,
    xleftmargin=1em,
    xrightmargin=1em,
    aboveskip=1em,
    belowskip=1em,
    keywordstyle=\color{codekeyword}\bfseries,
    commentstyle=\color{codecomment}\itshape,
    stringstyle=\color{codestring},
}

% YAML language definition for listings
\lstdefinelanguage{YAML}{
    keywords={true,false,null,y,n},
    keywordstyle=\color{blue}\bfseries,
    sensitive=false,
    comment=[l]{\#},
    morecomment=[s]{/*}{*/},
    commentstyle=\color{gray}\ttfamily,
    stringstyle=\color{red}\ttfamily,
    morestring=[b]',
    morestring=[b]",
}

% TOML language definition for listings
\lstdefinelanguage{TOML}{
    keywords={true,false},
    keywordstyle=\color{blue}\bfseries,
    sensitive=true,
    comment=[l]{\#},
    commentstyle=\color{gray}\ttfamily,
    stringstyle=\color{red}\ttfamily,
    morestring=[b]',
    morestring=[b]",
}

% Dockerfile language definition for listings
\lstdefinelanguage{Dockerfile}{
    keywords={FROM,RUN,CMD,LABEL,EXPOSE,ENV,ADD,COPY,ENTRYPOINT,VOLUME,USER,WORKDIR,ARG,ONBUILD,STOPSIGNAL,HEALTHCHECK,SHELL},
    keywordstyle=\color{blue}\bfseries,
    sensitive=false,
    comment=[l]{\#},
    commentstyle=\color{gray}\ttfamily,
    stringstyle=\color{red}\ttfamily,
    morestring=[b]',
    morestring=[b]",
}

% JSON language definition for listings
\lstdefinelanguage{JSON}{
    keywords={true,false,null},
    keywordstyle=\color{blue}\bfseries,
    sensitive=true,
    commentstyle=\color{gray}\ttfamily,
    stringstyle=\color{red}\ttfamily,
    morestring=[b]",
}

% Code listing float environment
\usepackage{float}
\newfloat{listing}{htbp}{lol}[section]
\floatname{listing}{Listing}

% Admonition boxes
\usepackage{tcolorbox}
\tcbuselibrary{skins,breakable}

% Admonition style definitions
\newtcolorbox{admonitionbox}[2][]{
    enhanced,
    breakable,
    colback=#2!5,
    colframe=#2!80!black,
    fonttitle=\bfseries,
    title=#1,
    left=2mm,
    right=2mm,
}

% Synthon tuple boxes
\newtcolorbox{synthonbox}{
    enhanced,
    colback=blue!5,
    colframe=blue!75!black,
    fontupper=\large,
    center upper,
    boxrule=1pt,
    arc=4pt,
}

\begin{document}

\title{SynthOmnicon: Typed Programming Language for Matter}
\date{\today}

\maketitle

\textbf{Status:} Phase 3a complete $\cdot$ Phase 3d complete (v0.4.0) $\cdot$ Phase 3b/3c pending $\cdot$ Phase 3e in progress (v0.4.4)

\begin{center}
\rule{0.5\textwidth}{0.4pt}
\end{center}

\subsection{Vision}

The Synthonicon is a \textbf{type system for matter} --- not just chemistry. Its eleven primitives are type constructors spanning molecular, supramolecular, temporal, quantum, and topological domains. Its axioms are refinement constraints. HotSwap transitions are category morphisms. Retrodesign is type inference. The Varma probe embeds proof obligations. The phase diagram maps the metric geometry of the type space itself.

v0.4.0 revealed something deeper than expected: the primitive algebra is \textbf{domain-agnostic by construction}. When quantum particles and topological materials were encoded using the same eleven primitives, the algebra immediately produced physically correct results --- non-Abelian anyons outranking ℤ-class invariants in tensor products, the K\_trap$\rightarrow$K\_MBL universal cost (+2.303 nats) appearing identically across three different topological phases, the Ward dendrogram separating extended topological matter from point particles at d$\approx$9.52 with no physics input. The framework did not simulate physics; it discovered it from the ordinal structure of the type system.

The goal of this track is to make that implicit structure \textbf{explicit and mechanically enforced} --- moving from ''axioms checked at runtime via string matching'' toward ''illegal states unrepresentable by construction,'' and from ''phase boundaries detected empirically'' toward ''phase boundaries provable from the Kleisli metric.''

\textbf{The primitives are relational operators, not intrinsic attributes.} Every element of the tuple describes a constraint between entities or a capacity for interaction. F (fidelity) is reliability of constraint satisfaction relative to a competitor --- there is no intrinsic F, only F relative to a context. K (kinetic) is a barrier to rearrangement, implying at least two states. $\Omega$ (topological protection) is protection against perturbations, meaningless without an environment. $\Gamma$ (grammar) is partner selection logic by definition. This is not a philosophical gloss; it is a \textbf{type-system requirement}: you cannot assign F, K, $\Gamma$, or $\Omega$ without specifying an interaction context. A synthon tuple without a context is a description of interaction potential, not isolated being.

The algebra enforces this. Every operation in \texttt{meet / join / tensor / path / lift / pipeline} requires at least one additional operand. \texttt{tensor(s₁, s₂)} computes mutual information between two systems. \texttt{path(src, dst)} requires two endpoints. \texttt{lift(s, critical)} requires a named target context and is blocked by a relational gate (\texttt{F \textgreater{}= $F_{\hbar}$}). There are no unary information generators. The algebra cannot process ''nothing but the object.'' This is the deeper reason the type system works: \textbf{it is a calculus of relations, and relations require at least two terms}.

Furthermore, the algebra is systematically asymmetric --- \texttt{path(A-\textgreater{}B) != path(B-\textgreater{}A)}, the F-floor ratchet is directed, \texttt{lift} has no inverse. This places the framework in the tradition of structural realism rather than symmetric relational theories: the world's causal structure is relational but ordered, and the ordering is the load-bearing part.

A design program in the finished language looks like:

\begin{lstlisting}[language=YAML]
# design.syn
version: "1.0"
start: soai_pyrimidyl_autocatalytic_cycle
do:
  - join: proline_aldol_cycle
  - lift: critical
  - path: varma_qxy_reference
    xi_tolerance: 1.5
  - assert:
      expr: phi_c_score > 0.70
      message: "Must approach Phi_c before ensemble step"
  - tensor: db24c8_dialkylammonium_pseudorotaxane
output:
  format: text
  save: result.json
\end{lstlisting}

\texttt{syncon run design.syn} evaluates this as a proper monadic do-block with typed effects, inline proof obligations, and full step tracing.

\begin{center}
\rule{0.5\textwidth}{0.4pt}
\end{center}

\subsection{Current State (v0.4.0)}

\begin{table}[htbp]
\centering
\small
\rowcolors{1}{white}{white}
\begin{tabularx}{\textwidth}{>{\centering\arraybackslash}X >{\centering\arraybackslash}X >{\centering\arraybackslash}X}
\toprule
\rowcolor{gray!25}\textbf{Component} & \textbf{Status} & \textbf{Notes} \\
\midrule
\rowcolors{1}{white}{gray!10}
Eleven-primitive tuple &  complete & \langleD;T;R;P;F;K;G;$\Gamma$;$\Phi$;S;$\Omega$\rangle --- $\Omega$ optional, defaults None for classical \\
Axioms 1--7 &  complete & Runtime; Axiom 6 discrete/continuous; Axiom 1 is quantum boundary detector \\
Algebra (meet/join/tensor/lift/path/distance) &  complete & \texttt{algebra.py}; all commands wired to \texttt{syncon} \\
DesignPipeline (Writer+Maybe) &  complete & Fluent interface + \texttt{bind}/\texttt{mzero}/\texttt{mplus} via SynthonM \\
F-floor HotSwap ratchet &  validated & 6/6 CB{[}7{]} experimental match \\
Varma probe (8 mechanisms) &  validated & Factors 1--5 QXY $\cdot$ Factor 6 steric-cliff $\cdot$ Factor 7 Frank $\cdot$ Factor 8 quantum criticality \\
\texttt{.syn} DSL runner &  complete & \texttt{syncon run design.syn} --- \texttt{syn\_runner.py} \textasciitilde{}380 lines \\
\texttt{SynthonM} monad type &  complete & \texttt{monad.py} --- WriterT+StateT+MaybeT stack \\
\texttt{DesignStrategy} type + combinators &  complete & \texttt{strategy\_then}, \texttt{strategy\_or}, \texttt{optimize} \\
Quantum primitive extensions &  complete & T\_braid $\cdot$ K\_MBL $\cdot$ $\Gamma$\_↓/QUANTUM $\cdot$ $\Omega$ (TopoIndex) --- Phase 3d \\
Quantum catalog (8 synthons) &  complete & \texttt{domains/quantum/} --- particles + topological matter, all $\Omega$ assigned \\
Tuple-space phase diagram &  complete & \texttt{phase\_diagram.py} --- Ward clustering + MDS + \texttt{syncon phase-diagram} \\
Refinement types (Pydantic v2) &  pending & Phase 3b target \\
\texttt{from\_disorder} KineticCharacter constructor &  pending & Phase 3b --- K\_MBL cannot use \texttt{from\_barrier} \\
\texttt{context.irreversible} in SynthonM &  pending & Phase 3b --- $\Gamma$\_↓ dissipative flag propagation \\
Kleisli category (formal) &  pending & Phase 3c target \\
Z3/SMT retrodesign &  pending & Phase 3c target \\
Lean 4 axiom formalization &  pending & Phase 3c long-term \\
Holographic dimensionality D\_holo &  complete & Phase 3e --- \texttt{Dimensionality.HOLOGRAPHIC}; \texttt{ads\_cft\_boundary} synthon registered \\
Phase transitions as morphisms &  complete & Phase 3e --- \texttt{synthomnicon/morphism.py}; \texttt{syncon transition} CLI; Kleisli arrow classification \\
Kitaev honeycomb B-phase synthon &  pending & Phase 3e --- T\_braid + $\Omega$\_NA + D\_△, 2D non-Abelian \\
MBL phase synthon (disordered Kitaev) &  pending & Phase 3e --- K\_MBL + T\_ \\
\bottomrule
\end{tabularx}
\end{table}

\begin{center}
\rule{0.5\textwidth}{0.4pt}
\end{center}

\subsection{Phase 3a --- Monadic Foundation}

\textbf{Target: \texttt{syncon run design.syn} works end-to-end.}

\subsubsection{3a.1 \texttt{synthomnicon/monad.py} --- SynthonM}

\texttt{SynthonM{[}A{]}} is the monad transformer stack:

\begin{lstlisting}
SynthonM[A] \cong WriterT[float] (StateT[Context] (MaybeT Identity)) A
\end{lstlisting}

Concretely, a \texttt{SynthonM{[}Synthon{]}} carries:

\begin{table}[htbp]
\centering
\small
\rowcolors{1}{white}{white}
\begin{tabularx}{\textwidth}{>{\centering\arraybackslash}X >{\centering\arraybackslash}X >{\centering\arraybackslash}X}
\toprule
\rowcolor{gray!25}\textbf{Effect} & \textbf{Carrier} & \textbf{Meaning} \\
\midrule
\rowcolors{1}{white}{gray!10}
\textbf{MaybeT} & \texttt{value: Optional{[}A{]}} & Computation may fail (BLOCKED/ERROR) \\
\textbf{WriterT} & \texttt{cost: float} & Accumulated $\Delta$\xi\_CP across all steps \\
\textbf{StateT} & \texttt{context: Context} & F-floor, criticality gate, step count \\
\textbf{Log} & \texttt{log: List{[}StepRecord{]}} & Full step trace for human inspection \\
\bottomrule
\end{tabularx}
\end{table}

\textbf{Monad primitives:}

\begin{lstlisting}[language=Python]
return_(a)        # Pure: lift value, zero cost
m.bind(f)         # Sequence: run m, feed into f, merge effects
m >> f            # Operator alias for bind
mzero()           # Failed computation
m1.mplus(m2)      # Try m1; if None, use m2
m1 | m2           # Operator alias for mplus
\end{lstlisting}

\textbf{Monadic lifts} wrap each algebra operation:

\begin{lstlisting}[language=Python]
join_m(name)            # algebra.join -> SynthonM
meet_m(name)            # algebra.meet -> SynthonM
tensor_m(name, lambda)       # algebra.tensor -> SynthonM
lift_m(target)          # _LIFT_MAP[target] -> SynthonM
path_m(name, \xi_tol)    # algebra.find_path -> SynthonM
assert_m(pred, msg)     # inline proof obligation -> SynthonM
\end{lstlisting}

\textbf{Strategy type:}

\begin{lstlisting}[language=Python]
DesignStrategy = Callable[[Synthon], SynthonM[Synthon]]

# Sequential composition
strategy_then(s1, s2)  # s1 >> s2

# Alternative / fallback
strategy_or(s1, s2)    # s1 <|> s2

# First-success search
optimize(synthon, [s1, s2, s3])  # asum over strategy list
\end{lstlisting}

\subsubsection{3a.2 \texttt{synthomnicon/syn\_runner.py} --- .syn DSL Evaluator}

Parses \texttt{.syn} YAML files and evaluates them as \texttt{SynthonM} pipelines.

\textbf{Step dispatch table:}

\begin{table}[htbp]
\centering
\small
\rowcolors{1}{white}{white}
\begin{tabularx}{\textwidth}{>{\centering\arraybackslash}X >{\centering\arraybackslash}X >{\centering\arraybackslash}X}
\toprule
\rowcolor{gray!25}\textbf{Step key} & \textbf{Monadic lift} & \textbf{Example} \\
\midrule
\rowcolors{1}{white}{gray!10}
\texttt{join} & \texttt{join\_m(name)} & \texttt{- join: proline\_aldol\_cycle} \\
\texttt{meet} & \texttt{meet\_m(name)} & \texttt{- meet: nitroso\_radical} \\
\texttt{tensor} & \texttt{tensor\_m(name, lambda)} & \texttt{- tensor: db24c8\textbackslash{}n  lambda: 0.4} \\
\texttt{lift} & \texttt{lift\_m(target)} & \texttt{- lift: critical} \\
\texttt{path} & \texttt{path\_m(name, \xi)} & \texttt{- path: varma\_qxy\textbackslash{}n  xi\_tolerance: 1.5} \\
\texttt{assert} & \texttt{assert\_m(pred)} & \texttt{- assert:\textbackslash{}n    expr: phi\_c\_score \textgreater{} 0.70} \\
\texttt{bind} & named strategy & \texttt{- bind: my\_strategy} \\
\bottomrule
\end{tabularx}
\end{table}

\textbf{\texttt{assert} expression grammar} (safe dispatch, no \texttt{eval}):

\begin{lstlisting}
phi_c_score > N        -> run varma probe, check score
phi_c_score >= N       -> same
fidelity == F_hbar     -> check synthon.fidelity
topology == T_bowtie   -> check synthon.topology
criticality_phase == Phi_c  -> check enum value
axiom6_satisfied       -> run AxiomValidator.validate_axiom6_temporal_grounding
gd_degeneracy == X     -> run probe, check gd_degeneracy_type
reset_type == discrete -> check grounding["reset"]["type"]
reset_type == continuous  -> same
\end{lstlisting}

\subsubsection{3a.3 CLI: \texttt{syncon run}}

\begin{lstlisting}[language=bash]
syncon run design.syn
syncon run design.syn --format json
syncon run design.syn --dry-run   # parse + validate without executing
syncon run design.syn --save result.json
\end{lstlisting}

\textbf{Output (text mode):}

\begin{lstlisting}
SynthOmnicon Run: design.syn
  Start: soai_pyrimidyl_autocatalytic_cycle

  1. [PASS]  join(proline_aldol_cycle)  Delta\xi=+0.000 nat
  2. [BLOCKED]  lift(critical)  --- F floor not met (F_eth < F_hbar required)
  3. [PASS]  assert(phi_c_score > 0.70)  --- ASSERT_FAIL: score=0.380

  Total Delta\xi_CP: +0.000 nat  |  Steps: 3  |  FAILED at step 2
\end{lstlisting}

\begin{center}
\rule{0.5\textwidth}{0.4pt}
\end{center}

\subsection{Phase 3b --- Refinement Types}

\textbf{Target: illegal states are unrepresentable, not just runtime-rejected.}

\subsubsection{3b.1 Pydantic v2 validators on Synthon}

Move all 7 axioms into \texttt{@model\_validator(mode='after')} decorators. Construction of an invalid Synthon raises \texttt{ValidationError}, not a \texttt{grounding\_failure} flag.

\begin{lstlisting}[language=Python]
class Synthon(BaseModel):
    fidelity: Fidelity
    topology: Topology
    # ...

    @model_validator(mode='after')
    def axiom1_fidelity_floor(self):
        if self.topology == Topology.CYCLIC_BOWTIE:
            if self.fidelity == Fidelity.LOW:
                raise ValueError("Axiom 1: T_⋈ requires F >= F_eth")
        return self
\end{lstlisting}

\subsubsection{3b.2 Ordered primitive types}

\begin{lstlisting}[language=Python]
class FidelityGe(Annotated[Fidelity, Ge(Fidelity.MEDIUM)]):
    """F >= F_eth --- satisfied by F_eth or F_hbar"""

CriticalityCandidateSynthon = Synthon[fidelity=FidelityGe]
\end{lstlisting}

Retrodesign returns \texttt{list{[}CriticalityCandidateSynthon{]}} for $\Phi$\_c targets --- caller knows statically that every leaf has F $\geq$ F\_eth.

\subsubsection{3b.3 Protocol classes for domains}

\begin{lstlisting}[language=Python]
class TemporalSynthon(Protocol):
    """Structural subtype requiring D_inf and grounding["reset"] block."""
    @property
    def dimensionality(self) -> Literal[Dimensionality.TEMPORAL]: ...
    @property
    def grounding(self) -> dict:
        assert "reset" in self.grounding
\end{lstlisting}

\texttt{trajectory validate} and \texttt{retrodesign} for D\_$\infty$ targets return \texttt{TemporalSynthon} so the type checker catches missing reset blocks before execution.

\begin{center}
\rule{0.5\textwidth}{0.4pt}
\end{center}

\subsection{Phase 3c --- Category Theory Embedding}

\textbf{Target: the algebra is provably correct and retrodesign is constraint-satisfaction.}

\subsubsection{3c.1 Explicit Kleisli category}

Document the category formally:

\begin{itemize}
    \item \textbf{Objects}: Synthons (equivalence classes of valid ten-tuples)
    \item \textbf{Morphisms}: HotSwap transitions \texttt{f: A -\textgreater{} B} with cost \texttt{Delta\xi\_CP(f) in ℝ\textgreater{}=0}
    \item \textbf{Identity}: trivial self-swap, cost 0
    \item \textbf{Composition}: \texttt{find\_path} (BFS composes morphisms; cost = sum of $\Delta$\xi\_CP per hop)
    \item \textbf{Enrichment}: over \texttt{(ℝ\textgreater{}=0, +, 0)} --- a Lawvere metric space
    \item \textbf{Monoidal product}: \texttt{tensor} (bifunctor; F-bottleneck as monoidal unit constraint)
    \item \textbf{Natural transformations}: \texttt{lift\_*} functions (lift\_to\_temporal, etc.)
\end{itemize}

Add \texttt{synthomnicon/category.py} with formal composition and identity checks.

\subsubsection{3c.2 Z3/SMT for retrodesign}

Replace BFS in \texttt{retrodesign.py} with Z3 constraint satisfaction:

\begin{lstlisting}[language=Python]
# target: \langleD_inf; T_⋈; R_‡; P_DA; $F_{\hbar}$; ...; Phi_sub; 1:1\rangle
# find all catalog entries satisfying axioms 1,2,4,6 and
# compositionally derivable from the target tuple
solver = z3.Solver()
# encode the ten-tuple as Z3 bitvectors
# encode axioms as Z3 constraints
# enumerate all satisfying assignments
\end{lstlisting}

This gives \textbf{complete} retrodesign (no false negatives from BFS depth limit) and opens constraint-solving for multi-target design.

\subsubsection{3c.3 Lean 4 axiom formalization (long-term)}

Formalize the seven axioms as Lean 4 theorems. The Python implementation becomes a ''trusted kernel'' whose correctness is guaranteed by the proof. Key lemmas:

\begin{itemize}
    \item \texttt{axiom1\_preserved\_by\_hotswap}: HotSwap preserves F $\geq$ F\_eth for T\_⋈ entries
    \item \texttt{join\_idempotent}: \texttt{join(s, s) = s}
    \item \texttt{meet\_join\_absorption}: \texttt{meet(s, join(s, t)) = s}
    \item \texttt{tensor\_fidelity\_bottleneck}: \texttt{tensor(s1, s2).F = min(s1.F, s2.F)}
\end{itemize}

\begin{center}
\rule{0.5\textwidth}{0.4pt}
\end{center}

\subsection{File Index}

\begin{table}[htbp]
\centering
\small
\rowcolors{1}{white}{white}
\begin{tabularx}{\textwidth}{>{\centering\arraybackslash}X >{\centering\arraybackslash}X >{\centering\arraybackslash}X}
\toprule
\rowcolor{gray!25}\textbf{File} & \textbf{Phase} & \textbf{Description} \\
\midrule
\rowcolors{1}{white}{gray!10}
\texttt{synthomnicon/monad.py} & 3a & SynthonM stack, monadic lifts, DesignStrategy \\
\texttt{synthomnicon/syn\_runner.py} & 3a & .syn YAML DSL evaluator \\
\texttt{synthomnicon/cli.py} & 3a/3d/3e & \texttt{syncon run} $\cdot$ algebra commands $\cdot$ \texttt{syncon phase-diagram} \\
\texttt{synthomnicon/models.py} & 3d & T\_braid $\cdot$ K\_MBL $\cdot$ QUANTUM/DISSIPATIVE grammar $\cdot$ TopoIndex \\
\texttt{synthomnicon/algebra.py} & 3d & $\Omega$ lattice ops $\cdot$ eleven-dimensional \texttt{tuple\_distance} \\
\texttt{synthomnicon/varma\_probe.py} & 3d & Factor 8 --- quantum criticality block \\
\texttt{synthomnicon/domains/quantum/\_\_init\_\_.py} & 3d & 8 canonical quantum/topological synthons \\
\texttt{synthomnicon/phase\_diagram.py} & 3e & Ward clustering $\cdot$ MDS $\cdot$ \texttt{PhaseDiagram} dataclass \\
\texttt{synthomnicon/morphism.py} & 3e & \texttt{TransitionMorphism} $\cdot$ \texttt{find\_transition()} $\cdot$ Kleisli arrow classification $\cdot$ \texttt{syncon transition} \\
\texttt{synthomnicon/category.py} & 3c & Kleisli category formal definition (pending) \\
\texttt{synthomnicon/typed.py} & 3b & Pydantic v2 Synthon, refinement types (pending) \\
\bottomrule
\end{tabularx}
\end{table}

\begin{center}
\rule{0.5\textwidth}{0.4pt}
\end{center}

\subsection{Design Decisions}

\textbf{Why WriterT + StateT + MaybeT (not ExceptT)?}
ExceptT short-circuits on the first failure and discards subsequent state. MaybeT propagates the accumulated cost and log even on failure, preserving the full trace. This matches the existing \texttt{DesignPipeline} semantics where BLOCKED steps are logged but execution records continue.

\textbf{Why YAML for \texttt{.syn} (not a custom syntax)?}
YAML is already parsed everywhere in the project; it's readable and diffable; LLMs generate it reliably. A custom syntax would require a parser that buys nothing concrete at this stage. The DSL can migrate to a typed syntax (like Dhall or a subset of Haskell do-notation) in Phase 3b once the semantics are stable.

\textbf{Why safe dispatch for \texttt{assert} expressions (not \texttt{eval})?}
\texttt{eval()} with arbitrary Python is a security surface and makes the language unpredictable. The predefined assertion grammar is a subset of useful predicates; any pattern not in the dispatch table raises a clear \texttt{UnknownAssertion} error pointing to the docs.

\textbf{Why not rewrite in Haskell?}
The Python catalog, CLI, and grounding layer are the validated, empirically-anchored core. Rewriting loses that incrementally. The right path is to formalize the semantics in Python first (Phase 3a/3b), then extract the formal core to Lean 4 (Phase 3c) as a proof of correctness --- not a replacement.

\begin{center}
\rule{0.5\textwidth}{0.4pt}
\end{center}

\begin{center}
\rule{0.5\textwidth}{0.4pt}
\end{center}

\subsection{Phase 3d --- Quantum Primitive Extensions (v0.4.0)}

\textbf{Motivation.} Encoding five quantum particles (photon, proton, electron, spin, qubit) through the LLM agent using only the classical primitive set revealed seven structural gaps --- places where the framework failed to encode known physics, and where each failure pointed at a real but undescribed class of physical system. The gaps were filled by adding four new primitive values and one new primitive.

\subsubsection{New primitive values}

\begin{table}[htbp]
\centering
\small
\rowcolors{1}{white}{white}
\begin{tabularx}{\textwidth}{>{\centering\arraybackslash}X >{\centering\arraybackslash}X >{\centering\arraybackslash}X >{\centering\arraybackslash}X}
\toprule
\rowcolor{gray!25}\textbf{Primitive} & \textbf{New value} & \textbf{Encodes} & \textbf{Physical systems} \\
\midrule
\rowcolors{1}{white}{gray!10}
T & \texttt{T\_braid} & Anyonic/braided exchange statistics & Fractional QHE ($\nu$=1/3, $\nu$=5/2), Kitaev honeycomb B-phase \\
K & \texttt{K\_MBL} & Many-body localization --- disorder-frozen & Disordered quantum magnets, Aubry-André cold atoms \\
$\Gamma$ operator & \texttt{DISSIPATIVE} ($\Gamma$\_↓) & Irreversible information loss & Lindblad open systems, quantum Zeno effect \\
$\Gamma$ tier & \texttt{QUANTUM} & Superposition-preserving (Toffoli semantics) & Quantum AND-gates, fault-tolerant circuits \\
\bottomrule
\end{tabularx}
\end{table}

\textbf{K\_MBL lattice position:} ordinal 0, below K\_trap (ordinal 1). MBL is more kinetically arrested than any energy-barrier trap --- the system cannot relax even with unlimited energy input (the many-body eigenbasis structure prevents thermalization). \texttt{meet(K\_MBL, K\_fast) -\textgreater{} K\_MBL}; \texttt{tensor(K\_MBL, anything) -\textgreater{} K\_MBL}.

\textbf{T\_braid tensor rule:} \texttt{tensor(T\_braid, T\_braid) -\textgreater{} T\_braid}. Two braided systems form a larger braided system; anyonic statistics do not network-promote like spatial structures. \texttt{T\_braid ⊓ T\_linear = \perp} (the unresolvable conflict that motivated the primitive).

\subsubsection{New primitive $\Omega$ (11th)}

\texttt{TopoIndex} encodes the symmetry class of topological protection, based on the Altland--Zirnbauer / K-theory (10-fold way) classification:

\begin{lstlisting}
Omega_0  (TRIVIAL)     --- No topological protection; classical systems
Omega_Z  (Z_CLASS)     --- ℤ winding number: Kitaev chain, SSH model, 1D p-wave
Omega_Z2 (Z2_CLASS)    --- ℤ₂ invariant: HgTe/CdTe QWs, Bi₂Se₃ (class AII/DIII)
Omega_C  (CHERN)       --- Integer Chern number: IQH, Chern insulators (class A)
Omega_NA (NON_ABELIAN) --- Non-abelian anyons: nu=5/2 FQH, Kitaev honeycomb
\end{lstlisting}

Protection-strength ordinal: TRIVIAL(0) \textless{} Z2(1) \textless{} Z(2) \textless{} CHERN(3) \textless{} NON\_ABELIAN(4).

In \texttt{meet}: lower protection propagates (conservative guarantee).
In \texttt{join} and \texttt{tensor}: higher protection propagates (strongest class dominates).
In \texttt{tuple\_distance}: \texttt{Omega} weight = 0.7 (highest categorical penalty --- a Chern insulator and a trivial metal are categorically different systems).

\subsubsection{Factor 8 --- Quantum criticality in the Varma probe}

The spin singlet analysis showed that the Varma probe scored 0.0 for spin despite \texttt{G\_aleph}, because all five classical factors require \texttt{D\_inf}. Factor 8 fires on the orthogonal pattern:

\begin{lstlisting}
G_aleph + F_hbar + K_trap + notD_inf  ->  quantum criticality (TFI/heavy_fermion class)
\end{lstlisting}

Weight 0.20. Falsifiable prediction: $\chi(T \to 0) \sim T^{-\gamma}$. The spin singlet scores 0.20 with universality class ''quantum\_criticality (TFI/heavy\_fermion)'', triggering \texttt{gd\_degeneracy\_type = ''quantum\_ground\_state''} and Axiom 5 satisfaction.

\subsubsection{Type-theoretic implications}

The new primitives expose three open questions for Phase 3b:

\begin{enumerate}
    \item \textbf{$\Omega$ is the first non-enum primitive candidate}: it has a natural integer (protection strength) representation in addition to its categorical (symmetry class) representation. Phase 3b refinement types could encode \texttt{Omega\_NA \textgreater{}= Omega\_Z} as a Pydantic \texttt{Ge} constraint.
    \item \textbf{K\_MBL violates \texttt{from\_barrier} semantics}: disorder-induced localization cannot be assigned from a single $\Delta G^{\ddagger}$ value. The \texttt{from\_barrier} classmethod correctly omits K\_MBL, but a new \texttt{from\_disorder} classmethod is implied.
    \item \textbf{$\Gamma$\_↓(DISSIPATIVE) introduces irreversibility into the grammar}: the existing \texttt{SynthonM} monad tracks cost accumulation (WriterT) but has no irreversibility primitive. A dissipative synthon flowing through a \texttt{SynthonM} pipeline should perhaps trigger a \texttt{StateT} flag: \texttt{context.irreversible = True}.
\end{enumerate}

\begin{center}
\rule{0.5\textwidth}{0.4pt}
\end{center}

\subsection{Phase 3e --- Tuple-Space Phase Detection (v0.4.0)}

\textbf{Target: the framework discovers its own phase boundaries from the metric geometry of the type space.}

\subsubsection{3e.1 \texttt{synthomnicon/phase\_diagram.py} --- PhaseDiagram module}

\texttt{PhaseDiagram} is a dataclass computed from the full N$\times$N pairwise \texttt{tuple\_distance} matrix over a named set of synthons:

\begin{lstlisting}[language=Python]
@dataclass
class PhaseDiagram:
    synthon_names: List[str]
    distance_matrix: np.ndarray       # NxN symmetric
    candidates: List[PhaseBoundary]   # ranked boundary list
    linkage_matrix: np.ndarray        # Ward linkage (scipy format)
    mds_coords: np.ndarray            # Nx2 MDS embedding
    factor8_flags: Dict[str, bool]    # Factor-8 trigger per synthon
    omega_values: Dict[str, str]      # \Omega class per synthon
    k_trap_flags: Dict[str, bool]     # K_trap flag per synthon
\end{lstlisting}

\texttt{build\_phase\_map(synthon\_names=None, catalog=None) -\textgreater{} PhaseDiagram} constructs the full object from the catalog in one call.

\textbf{Phase boundary detection algorithm:}

\begin{enumerate}
    \item Compute N$\times$N distance matrix via \texttt{tuple\_distance}.
    \item Run Ward hierarchical clustering on the upper triangle.
    \item Extract consecutive height differences in the merge sequence.
    \item Rank gaps as Major (top 33\%), Intermediate (middle), Minor (bottom 33\%).
    \item The primary boundary (largest gap, d $\approx$ 9.52 for the 8-synthon quantum catalog) is the most semantically significant cluster split.
\end{enumerate}

\begin{lstlisting}[language=Python]
from synthomnicon.phase_diagram import build_phase_map
pd = build_phase_map()
pd.print_report()
pd.plot(save_path="diagram.png")
\end{lstlisting}

\subsubsection{3e.2 \texttt{syncon phase-diagram} CLI command}

\begin{lstlisting}[language=bash]
# All synthons in catalog:
syncon phase-diagram

# Subset by name:
syncon phase-diagram photon proton electron kitaev_chain_majorana

# Save rendered figure:
syncon phase-diagram --save diagram.png

# Skip matplotlib (text + JSON only):
syncon phase-diagram --text-only

# JSON output for scripting:
syncon phase-diagram --format json
\end{lstlisting}

Output (text mode): distance matrix $\cdot$ phase boundary table (gap, label, threshold d) $\cdot$ Factor-8 cluster summary $\cdot$ K\_trap/$\Omega$ membership list.

\subsubsection{3e.3 MDS embedding --- metric geometry of the type space}

The 2-D MDS projection uses classical (metric) MDS via eigendecomposition of the double-centred Gram matrix:

\begin{lstlisting}
B = -½ H D^2 H        where H = I - (1/N)𝟏𝟏ᵀ
B ~= V₂ \Lambda₂ V₂ᵀ       (top-2 eigenpairs)
X = V₂ \Lambda₂^{1/2}     (Nx2 coordinates)
\end{lstlisting}

No sklearn dependency --- pure \texttt{numpy.linalg.eigh}. Stress is not reported; the Ward dendrogram is the primary clustering artifact.

\textbf{Key observation from the 8-synthon quantum map:}

\begin{itemize}
    \item 
    \begin{itemize}
    \item The primary boundary at d $\approx$ 9.52 separates two branches with no physics input:\textbf{Branch A (extended topological):} kitaev\_chain\_majorana $\cdot$ fqh\_moore\_read $\cdot$ topological\_insulator\_bi2se3
    \item \textbf{Branch B (point particles):} photon $\cdot$ proton $\cdot$ electron $\cdot$ spin\_singlet $\cdot$ qubit\_logical
\end{itemize}

    \item Within Branch A, fqh\_moore\_read (T\_braid + $\Omega$\_NA) separates from the ℤ-class pair at d $\approx$ 4.3.
    \item Within Branch B, proton and electron cluster at d $\approx$ 1.8 (charge sign only); qubit separates last at d $\approx$ 3.2.
\end{itemize}

The framework isolated the topological phase boundary by ordinal arithmetic alone --- no Hamiltonian, no symmetry group, no physical input beyond primitive assignments.

\subsubsection{3e.4 Universality track prediction}

The perturbation sweep \texttt{K\_trap -\textgreater{} K\_MBL} produces $\Delta$\xi = +2.303 nats {[}HIGH{]} for spin\_singlet, kitaev\_chain\_majorana, and fqh\_moore\_read --- three synthons with different T and $\Omega$ values, all producing the same thermodynamic cost. This is a \textbf{framework prediction}: the cost of entering MBL from a coherent gap-protected state is independent of the specific topological invariant.

The MDS map should display this as three parallel trajectories (same displacement vector, different starting positions) --- this is the ''universality track'' that Phase 3e will make visually and algebraically explicit.

\textbf{Falsification path:} thermodynamic integration or quench spectroscopy on Kitaev-chain vs FQH vs spin-singlet systems should show the same entropy cost to disorder-freeze the state.

\textbf{Contrasting result:} topological\_insulator\_bi2se3, \texttt{K\_slow -\textgreater{} K\_moderate} = $-$0.847 nats {[}MEDIUM{]}. Improving surface mobility \emph{reduces} thermodynamic cost --- gapless surface states, not a gapped condensate. The sign flip is structural (K ordinal direction reverses semantics for metals vs insulators).

\subsubsection{3e.5 Phase 3e stress test agenda}

Three classes of objects identified by GPT-4 review as stress tests for the tuple-space phase detection:

\textbf{Phase transitions as morphisms  complete (v0.4.4)}
\texttt{synthomnicon/morphism.py} implements \texttt{TransitionMorphism} and \texttt{find\_transition(src, dst, catalog)}. A transition A$\rightarrow$B is classified as:

\begin{itemize}
    \item \textbf{2nd order} --- HotSwap path exists through $\Phi$\_c intermediates (continuous; \texttt{syncon transition} finds the path and reports forward/reverse costs and $\Phi$\_c intermediates)
    \item \textbf{1st order} --- No path (D/T structural conflict); virtual Kleisli arrow with infinite primitive cost; latent heat $\approx$ barrier height between D/T classes
\end{itemize}

Key asymmetry result: \texttt{topological\_insulator\_bi2se3 -\textgreater{} synthon\_Fermi\_liquid} is an asymmetric 1st-order morphism --- the reverse path (Fermi liquid $\rightarrow$ TI, fidelity \emph{increases}) is permitted; the forward path (TI $\rightarrow$ Fermi liquid, fidelity \emph{decreases}) is blocked by the F-floor. Topological protection encodes as morphism irreversibility. \texttt{syncon transition SRC DST} CLI registered.

\textbf{Floquet synthons (periodic drive)  complete (v0.4.4)}
\texttt{floquet\_chern\_insulator} ($D_{\bigtriangleup}$, $T_{\uparrow\downarrow}$, $K_{\text{trap}}$, $\Omega_C$, $\Gamma_{\to}(\text{SEQUENTIAL})$) and \texttt{time\_crystal\_dtc} ($D_{\infty}$, $T_{\bowtie}$, $K_{\text{MBL}}$, $\Omega_Z$) registered. $d(\text{floquet\_chern}, \text{time\_crystal}) = 4.2$ --- same $G_{\gimel}/F_{\eth}/\Phi_{\text{sub}}$ floor, differing in $D/T/K$ and topological class ($\Omega_C$ vs. $\Omega_Z$). $\Gamma_{\to}(\text{SEQUENTIAL})$ encodes the stroboscopic Floquet operator; $K_{\text{MBL}}$ encodes the disorder-localization required for DTC stability against Floquet heating.

\textbf{Gauge theories (loop variables)  complete (v0.4.4)}
\texttt{z2\_lattice\_gauge\_toric\_code} registered: $\langle D_{\bigtriangleup}; T_{\uparrow\downarrow}; R_{\Leftrightarrow}; P_{\pm}^{\text{sym}}; F_{\hbar}; K_{\text{trap}}; G_{\aleph}; \Gamma_{\wedge}(\text{QUANTUM}); \Phi_{\text{sub}}; \Omega_{Z_2} \rangle$. Key result: $d(\text{z2\_toric}, \text{fqh\_moore\_read}) = 3.90$ at path cost = \textbf{0.000 nat}. Zero-cost path at non-zero tuple distance = same thermodynamic universality class, different topological invariant. $T_{\uparrow\downarrow}$ (Wilson loop holonomy), $R_{\Leftrightarrow}$ (loop variable recognition), $G_{\aleph}$ (global stabilisers) encode the gauge-theory structure without a new primitive.

\begin{center}
\rule{0.5\textwidth}{0.4pt}
\end{center}

\subsection{Algebraic Operations Reference}

\emph{Runnable companion: \texttt{TENSOR\_OPS\_DEMO.py} --- 18 worked examples, \texttt{python TENSOR\_OPS\_DEMO.py --section \textless{}op\textgreater{}}.}

The synthon tuple space is a product of bounded lattices. Seven operations act on it. Each maps cleanly to a concept from tensor mathematics, category theory, or constructive logic --- not as a metaphor but as a precise structural translation.

\subsubsection{Operation Table}

\begin{table}[htbp]
\centering
\small
\rowcolors{1}{white}{white}
\begin{tabularx}{\textwidth}{>{\centering\arraybackslash}X >{\centering\arraybackslash}X >{\centering\arraybackslash}X >{\centering\arraybackslash}X}
\toprule
\rowcolor{gray!25}\textbf{Operation} & \textbf{Signature} & \textbf{Category-theoretic reading} & \textbf{SynthonM lift} \\
\midrule
\rowcolors{1}{white}{gray!10}
\texttt{meet(s₁, s₂)} & \texttt{Synthon\^{}2 -\textgreater{} Synthon ⊔ \{\perp\}} & Categorical product; infimum in product lattice & \texttt{meet\_m(name)} \\
\texttt{join(s₁, s₂)} & \texttt{Synthon\^{}2 -\textgreater{} Synthon ⊔ \{\perp\}} & Coproduct; supremum; F-floor raised in WriterT & \texttt{join\_m(name)} \\
\texttt{tensor(s₁, s₂, lambda)} & \texttt{Synthon\^{}2 x {[}0,1{]} -\textgreater{} Synthon} & Bifunctor; \xi\_tensor = \xi₁ + \xi₂ $-$ $\lambda$$\cdot$I(s₁,s₂) & \texttt{tensor\_m(name, lambda)} \\
\texttt{lift(s, target)} & \texttt{Synthon -\textgreater{} Synthon} & Functor between domain categories & \texttt{lift\_m(target)} \\
\texttt{path(src, dst)} & \texttt{Synthon\^{}2 -\textgreater{} PathResult} & Geodesic in the Kleisli enriched category & \texttt{path\_m(name, \xi\_tol)} \\
\texttt{assert(pred)} & \texttt{(Synthon -\textgreater{} bool) -\textgreater{} SynthonM} & Inline proof obligation; does not alter value & \texttt{assert\_m(pred, msg)} \\
\textbf{decomp} & see below & Inversion, projection, factorisation & \texttt{cofactor\_m}, \texttt{factor\_m}, \texttt{project\_m} \\
\bottomrule
\end{tabularx}
\end{table}

\begin{center}
\rule{0.5\textwidth}{0.4pt}
\end{center}

\subsubsection{§1 $\cdot$ meet (⊓) --- Greatest Lower Bound}

\textbf{Primitive rules:}

\begin{itemize}
    \item \emph{Categorical} (D, T, R, P, $\Gamma$): identity required; mismatch $\rightarrow$ \perp
    \item \emph{Ordered} (F, K, G, $\Omega$): min --- most conservative
    \item \emph{$\Phi$}: \textbf{$\Phi$\_c is an absorbing element} --- $\Phi$\_c ⊓ $\Phi$\_sub = $\Phi$\_c
\end{itemize}

The absorbing $\Phi$\_c is a co-Heyting top: not the usual infimum but a constraint forcing the result into the critical sub-lattice. The algebra cannot ''average away'' a phase transition.

\textbf{Key examples (§1 in demo):}

\begin{table}[htbp]
\centering
\small
\rowcolors{1}{white}{white}
\begin{tabularx}{\textwidth}{>{\centering\arraybackslash}X >{\centering\arraybackslash}X >{\centering\arraybackslash}X}
\toprule
\rowcolor{gray!25}\textbf{Input} & \textbf{Output} & \textbf{Insight} \\
\midrule
\rowcolors{1}{white}{gray!10}
\texttt{meet(Hv1\_human\_open, AtHv1\_primed)} & PASS $\cdot$ $\Phi$\_c propagates & Cross-species water-chain conservation proved algebraically \\
\texttt{meet(Hv1\_human\_open, 2GBI\_inhibitor)} & \perp $\cdot$ T-conflict + P-conflict & Conflict \emph{is} the answer: inhibitor occludes, does not merge \\
\texttt{meet(cooper\_pair, topological\_insulator)} & PASS on $\Omega$, \perp on D/P/$\Gamma$ & $\Omega$\_Z ⊓ $\Omega$\_Z₂ = $\Omega$\_Z₂ (AZ classification ordinal) \\
\bottomrule
\end{tabularx}
\end{table}

\begin{center}
\rule{0.5\textwidth}{0.4pt}
\end{center}

\subsubsection{§2 $\cdot$ join (⊔) --- Least Upper Bound / Design Target}

\textbf{Primitive rules:}

\begin{itemize}
    \item \emph{Categorical}: identity required; mismatch $\rightarrow$ \perp
    \item \emph{Ordered} (F, K, G, $\Omega$): max --- most demanding; F-floor ratchet raised
    \item \emph{$\Phi$}: \textbf{$\Phi$\_c is join-dominant} --- $\Phi$\_c ⊔ $\Phi$\_sub = $\Phi$\_c
\end{itemize}

join is the \textbf{design target}: the minimal synthon that both inputs can inject into. In module theory: the pushout. The F-floor ratchet implements the \emph{directed} nature of commitment --- once a fidelity obligation is raised, it cannot be lowered by any downstream operation.

\textbf{Key examples (§2 in demo):}

\begin{table}[htbp]
\centering
\small
\rowcolors{1}{white}{white}
\begin{tabularx}{\textwidth}{>{\centering\arraybackslash}X >{\centering\arraybackslash}X >{\centering\arraybackslash}X}
\toprule
\rowcolor{gray!25}\textbf{Input} & \textbf{Output} & \textbf{Insight} \\
\midrule
\rowcolors{1}{white}{gray!10}
\texttt{join(Hv1\_human\_open, PsHv1\_constitutive)} & PASS $\cdot$ $\Phi$\_c propagates, F$\rightarrow$$F_{\hbar}$ & Gymnosperm channel is algebraically subsumed \\
\texttt{join(2GBI\_inhibitor, HIF\_inhibitor)} & \perp $\cdot$ T-conflict & No common scaffold --- the incompatibility is categorical, not quantitative \\
\texttt{join(imatinib, GNF-2)} & \perp $\cdot$ T-conflict reveals bottleneck & Dual-mechanism ABL design: T\_⋈ \cap T\_branched = \emptyset \\
\bottomrule
\end{tabularx}
\end{table}

\begin{center}
\rule{0.5\textwidth}{0.4pt}
\end{center}

\subsubsection{§3 $\cdot$ tensor (\otimes) --- Bifunctor / Co-Assembly Prediction}

\textbf{\xi cost formula:}


\[\xi_\text{tensor} = \xi_1 + \xi_2 - \lambda \cdot I(s_1, s_2)\]


where $I(s_1,s_2) \approx \frac{\text{primitive matches}}{7} \times \min(\xi_1, \xi_2)$ and $\lambda \in [0,1]$ is the mutual-information discount (shared structure reduces assembly cost).

\textbf{Primitive rules:}

\begin{itemize}
    \item D: union of domain sets
    \item T: \textbf{topology promotion} (cage \textgreater{} network \textgreater{} hub \textgreater{} bowtie \textgreater{} linear); same $\rightarrow$ same; \texttt{T\_braid \otimes T\_braid -\textgreater{} T\_braid}
    \item F, K: min (bottleneck propagates)
    \item G: max (coarsest scale controls)
    \item $\Phi$, $\Omega$: join-dominant (criticality and topological protection propagate)
\end{itemize}

\textbf{Critical semantic boundary:}

\begin{quote}
tensor predicts \textbf{co-occupancy} (the two-particle Hilbert space). A \textbf{bound state} requires \texttt{tensor} \emph{then} \texttt{meet(binding\_potential\_synthon)} to acquire T\_bowtie. This is the exciton theorem: \texttt{electron \otimes hole -\textgreater{} T\_\textbar{} \otimes T\_\textbar{} = T\_\textbar{}}; the Coulomb interaction that produces T\_⋈ is a \emph{separate} meet step.

\end{quote}

\textbf{Key examples (§3 in demo):}

\begin{table}[htbp]
\centering
\small
\rowcolors{1}{white}{white}
\begin{tabularx}{\textwidth}{>{\centering\arraybackslash}X >{\centering\arraybackslash}X >{\centering\arraybackslash}X}
\toprule
\rowcolor{gray!25}\textbf{Input} & \textbf{Output} & \textbf{Insight} \\
\midrule
\rowcolors{1}{white}{gray!10}
\texttt{tensor(electron, hole, lambda=0.5)} & PASS $\cdot$ T stays LINEAR $\cdot$ P$\rightarrow$DONOR\_ACCEPTOR & bound state $\neq$ tensor product; Coulomb binding is a meet \\
\texttt{tensor(phonon\_acoustic, magnon, lambda=0.4)} & PASS $\cdot$ F$\rightarrow$$F_{\ell}$ $\cdot$ G$\rightarrow$G\_ℵ & magnetoelastic polaron; $\lambda$ encodes spin-phonon coupling g\_kq \\
\texttt{tensor(majorana, majorana, lambda=0.3)} & PASS $\cdot$ T\_braid preserved $\cdot$ $\Omega$\_NA preserved & topological qubit: braid statistics do not network-promote \\
\bottomrule
\end{tabularx}
\end{table}

\textbf{Special T\_braid rule:} \texttt{T\_braid \otimes T\_braid -\textgreater{} T\_braid}. The braid group $B_n \otimes B_m \subseteq B_{n+m}$ --- anyonic statistics compose into a larger braid system, never into a network.

\begin{center}
\rule{0.5\textwidth}{0.4pt}
\end{center}

\subsubsection{§4 $\cdot$ lift --- Natural Transformations Between Domain Categories}

Three lifts are implemented as functors between domain categories:

\begin{table}[htbp]
\centering
\small
\rowcolors{1}{white}{white}
\begin{tabularx}{\textwidth}{>{\centering\arraybackslash}X >{\centering\arraybackslash}X >{\centering\arraybackslash}X >{\centering\arraybackslash}X}
\toprule
\rowcolor{gray!25}\textbf{Lift} & \textbf{Source $\rightarrow$ Target} & \textbf{Key primitive changes} & \textbf{Cost} \\
\midrule
\rowcolors{1}{white}{gray!10}
\texttt{lift\_to\_temporal} & C\_mol $\rightarrow$ C\_temporal & D$\rightarrow$D\_$\infty$, T$\rightarrow$T\_⋈, R$\rightarrow$R\_‡, $\Gamma$$\rightarrow$$\Gamma$\_$\rightarrow$(SEL) & 0.0 nats \\
\texttt{lift\_to\_spatial} & C\_mol $\rightarrow$ C\_crystal & D$\rightarrow$D\_△, T$\rightarrow$T\_□, G escalates to MESOSCALE & 0.0 nats \\
\texttt{criticality\_lift} & C\_sub $\rightarrow$ C\_critical & $\Phi$\_sub$\rightarrow$$\Phi$\_c & \textbf{+2.303 nats} \\
\bottomrule
\end{tabularx}
\end{table}

\textbf{Criticality lift cost = 2.303 nats} is the Shannon entropy of a binary phase transition: $\log_e(10)$, one decade of probability mass transfer. It is the primitive-space Landauer bound: the minimum cost of acquiring an order parameter.

\textbf{Gate:} \texttt{criticality\_lift} requires F $\geq$ $F_{\hbar}$. Blocked if F \textless{} $F_{\hbar}$. This is the relational gate that prevents phonons ($F_{\ell}$) from being criticality-lifted --- they are not the order parameter.

\textbf{Asymmetry:} lift has no inverse. The F-floor ratchet in SynthonM is monotone: \texttt{join} and \texttt{lift} raise the floor; no operation lowers it. This encodes thermodynamic irreversibility as a monad law.

\textbf{Category-theoretic translations:}

\begin{itemize}
    \item \texttt{lift\_to\_temporal} \cong loop-space functor $\Omega$: sends a pointed space (binding event) to its loop space (catalytic cycle), endowing it with a monoid structure
    \item \texttt{lift\_to\_spatial} \cong classifying space functor B(G): the SBU is the classifying object for the crystal's packing symmetry group
    \item \texttt{criticality\_lift} \cong the ''phase enrichment'' functor from a generic category to a $\Phi$\_c-structured category at an RG fixed point
\end{itemize}

\begin{center}
\rule{0.5\textwidth}{0.4pt}
\end{center}

\subsubsection{§5 $\cdot$ path --- Geodesic in the HotSwap Kleisli Category}

\textbf{Category:}

\begin{itemize}
    \item Objects: synthons (equivalence classes of 10-tuples)
    \item Morphisms: HotSwap transitions $f: A \to B$ with cost $\Delta\xi_\text{CP}(f) \geq 0$
    \item Enrichment: over $(\mathbb{R}_{\geq 0}, +, 0)$ --- a \textbf{Lawvere metric space}
    \item Composition: BFS path; total cost additive
\end{itemize}

\textbf{Hard topological gates (path blockers):}

\begin{itemize}
    \item D-class must be invariant along the path
    \item T-class must be invariant along the path
\end{itemize}

A \textbf{blocked path} = a T- or D-class boundary = \textbf{1st-order-like transition} (discrete jump, latent cost \textgreater{} 0). A \textbf{found path} = 2nd-order-like transition (continuous deformation within the same homotopy class).

\textbf{Key examples (§5 in demo):}

\begin{table}[htbp]
\centering
\small
\rowcolors{1}{white}{white}
\begin{tabularx}{\textwidth}{>{\centering\arraybackslash}X >{\centering\arraybackslash}X >{\centering\arraybackslash}X >{\centering\arraybackslash}X}
\toprule
\rowcolor{gray!25}\textbf{Source $\rightarrow$ Target} & \textbf{Status} & \textbf{Reason} & \textbf{Physical reading} \\
\midrule
\rowcolors{1}{white}{gray!10}
\texttt{AtHv1\_silent -\textgreater{} AtHv1\_primed} & BLOCKED & T\_$\in$ $\neq$ T\_⋈ & Mechanical priming is a discrete topology jump, not continuous deformation \\
\texttt{2GBI\_inhibitor -\textgreater{} HIF\_inhibitor} & BLOCKED & T\_$\in$ $\neq$ T\_\textbar{} & Scaffold evolution d=1.5 nats; SAR cannot bridge this \\
\texttt{topological\_insulator -\textgreater{} electron} & BLOCKED & D + T mismatch & Topological gap = blocked path; bulk-boundary correspondence as monad law \\
\bottomrule
\end{tabularx}
\end{table}

\textbf{Asymmetric morphisms} (from Phase 3e): \texttt{path(TI -\textgreater{} Fermi liquid)} and \texttt{path(Fermi liquid -\textgreater{} TI)} are both blocked, but the F-floor ratchet means the reverse costs more. Topological protection encodes as \textbf{morphism irreversibility}.

\begin{center}
\rule{0.5\textwidth}{0.4pt}
\end{center}

\subsubsection{§6 $\cdot$ SynthonM --- Monad Transformer Stack}

\begin{lstlisting}
SynthonM[A] \cong  WriterT[ℝ>=0] (StateT[Context] (MaybeT Identity)) A
\end{lstlisting}

\begin{table}[htbp]
\centering
\small
\rowcolors{1}{white}{white}
\begin{tabularx}{\textwidth}{>{\centering\arraybackslash}X >{\centering\arraybackslash}X >{\centering\arraybackslash}X}
\toprule
\rowcolor{gray!25}\textbf{Layer} & \textbf{Carrier} & \textbf{Semantics} \\
\midrule
\rowcolors{1}{white}{gray!10}
MaybeT & \texttt{value: Optional{[}A{]}} & Failure propagation: BLOCKED short-circuits \\
WriterT & \texttt{cost: float} & Accumulated $\Delta$\xi\_CP across all steps \\
StateT & \texttt{context: Context} & F-floor ratchet; criticality gate; step count \\
\bottomrule
\end{tabularx}
\end{table}

\textbf{Monad laws in design terms:}

\begin{itemize}
    \item \texttt{return\_(s) \textgreater{}\textgreater{} f = f(s)} --- starting fresh has no cost
    \item \texttt{(m \textgreater{}\textgreater{} f) \textgreater{}\textgreater{} g = m \textgreater{}\textgreater{} (f \textgreater{}\textgreater{} g)} --- pipeline order is associative
    \item \texttt{mzero() \textgreater{}\textgreater{} f = mzero()} --- blocked computations stay blocked
\end{itemize}

\textbf{\texttt{mplus} (\textless{}\textbar{}\textgreater{}) --- fallback:} \texttt{BLOCKED \textless{}\textbar{}\textgreater{} PASS = PASS}. Models branching design paths. The F-floor from a failed branch does NOT transfer to the fallback --- the monad correctly isolates commitment state.

\textbf{Key examples (§6 in demo):}

\begin{table}[htbp]
\centering
\small
\rowcolors{1}{white}{white}
\begin{tabularx}{\textwidth}{>{\centering\arraybackslash}X >{\centering\arraybackslash}X >{\centering\arraybackslash}X}
\toprule
\rowcolor{gray!25}\textbf{Pipeline} & \textbf{Result} & \textbf{What it proves} \\
\midrule
\rowcolors{1}{white}{gray!10}
\texttt{Hv1\_open \textgreater{}\textgreater{} meet(AtHv1\_primed) \textgreater{}\textgreater{} join(PsHv1)} & PASS $\cdot$ $\Delta$\xi=0 & Cross-species conservation; $\Phi$\_c preserved through full pipeline \\
\texttt{strategy\_A.mplus(strategy\_B)} & strategy\_B & BLOCKED\oplusPASS=PASS; F-floor isolation \\
\texttt{cooper\_pair \textgreater{}\textgreater{} tensor(TI\_surface, lambda=0.3)} & PASS $\cdot$ $\Omega$\_Z $\cdot$ $\Phi$\_c & Proximity effect: topological protection and criticality propagate \\
\bottomrule
\end{tabularx}
\end{table}

\begin{center}
\rule{0.5\textwidth}{0.4pt}
\end{center}

\subsubsection{§7 $\cdot$ Decompositions}

Decomposition operations \textbf{invert or factor} the algebraic operations. They are the analytic tools; §§1--6 are the synthetic tools.

\begin{table}[htbp]
\centering
\small
\rowcolors{1}{white}{white}
\begin{tabularx}{\textwidth}{>{\centering\arraybackslash}X >{\centering\arraybackslash}X >{\centering\arraybackslash}X}
\toprule
\rowcolor{gray!25}\textbf{Operation} & \textbf{Category-theoretic reading} & \textbf{Use} \\
\midrule
\rowcolors{1}{white}{gray!10}
\texttt{factor(s)} & Decrement morphism: step one ordinal to its sub-object & ''What is the largest proper sub-synthon?'' \\
\texttt{cofactor(C, A)} & \textbf{Inverse bifunctor}: find B s.t. tensor(A, B) $\approx$ C & Reverse-engineer the unknown component of a co-assembly \\
\texttt{kernel(s, probe)} & \textbf{Kernel of a predicate-morphism}: largest sub-s annihilated by $\phi$ & ''What can I remove before $\Phi$\_c signal disappears?'' \\
\texttt{principal\_decomp(s)} & \textbf{Join-irreducible basis decomposition} (Birkhoff representation) & SVD analog: ordered list of primitive contributions \\
\texttt{project(s, dims)} & \textbf{Coordinate projection} $\pi$\_S: retain S, zero out S̄ & Subspace restriction to named primitives \\
\texttt{complement\_rel(s, ctx, tgt)} & \textbf{Relative pseudocomplement} in a Heyting algebra & ''What part of s is still needed after ctx already covers tgt?'' \\
\texttt{primitive\_peel(s, prim)} & Single-dim projection with invariant cost accounting & Remove one primitive; pay $\Phi$\_c/$\Omega$ protection costs \\
\bottomrule
\end{tabularx}
\end{table}

\textbf{Key examples (§7 in demo):}

\textbf{cofactor} --- Cooper pair reverse engineering:

\begin{lstlisting}
cofactor(cooper_pair, conducting_electron) -> inferred partner B:
  K: BOTTLENECK   -> B must be K_slow (condensate timescale, not Fermi velocity)
  G: CONTRIBUTOR  -> B must be G_meso (coherence length \xi ~ 100--1000 nm)
  T: PASSTHROUGH  -> B must have T_⋈ (pairing loop)
  Phi: CONTRIBUTOR  -> B carries Phi_c (superfluid phase transition)
  \Omega: CONTRIBUTOR  -> B carries \Omega_Z (ℤ winding number)
\end{lstlisting}

The cofactor reconstructs the phonon-dressed retarded partner from the observable pair. This is how Cooper pair formation is diagnosed from first principles in primitive space.

\textbf{principal\_decomp} --- SVD of GNF-2:

\begin{lstlisting}
principal_decomp(GNF-2) -> 4 join-irreducible atoms:
  [1] atom[F=F_ℇ]     --- fidelity contribution    (most constraining)
  [2] atom[K=K_mod]   --- kinetic contribution
  [3] atom[G=G_ג]     --- mesoscale contribution
  [4] skeleton(GNF-2) --- categorical primitives (Phi_c, T_branched, ...)
\end{lstlisting}

Order = hardest to engineer first. If improving GNF-2, start at {[}3{]} (G escalation to mesoscale is the mechanistic bottleneck for allosteric propagation).

\textbf{project + complement\_rel} --- Heyting pseudocomplement:

\begin{lstlisting}[language=Python]
# "If the Phi/\Omega structure is already given by the topological material,
#  what does GNF-2 uniquely contribute toward the cooper_pair target?"
proj = project(cooper_pair, ["Phi", "Omega"])          # pi_{Phi,\Omega}(cooper_pair)
crel = complement_rel(gnf2, context=proj, target=cooper_pair)
# -> crel: K_slow + G_meso + T_branched (the part context does NOT cover)
\end{lstlisting}

The Heyting implication $\text{GNF-2} \Rightarrow \text{cooper\_pair}$ relative to $\neg\text{context}$: the maximal sub-synthon of GNF-2 that (a) has no overlap with the $\Phi$/$\Omega$ projection and (b) together with it covers the target. This is \textbf{design as constructive proof}.

\begin{center}
\rule{0.5\textwidth}{0.4pt}
\end{center}

\subsubsection{Canonical Pedagogical Example}

The Webster/Tombola Hv1 trilogy (Papers 1--3) threads through every operation:

\begin{table}[htbp]
\centering
\small
\rowcolors{1}{white}{white}
\begin{tabularx}{\textwidth}{>{\centering\arraybackslash}X >{\centering\arraybackslash}X >{\centering\arraybackslash}X}
\toprule
\rowcolor{gray!25}\textbf{Operation} & \textbf{Example} & \textbf{Result} \\
\midrule
\rowcolors{1}{white}{gray!10}
meet & \texttt{Hv1\_open ⊓ AtHv1\_primed} & $\Phi$\_c preserved; cross-species conservation \\
join & \texttt{Hv1\_open ⊔ PsHv1} & Gymnosperm subsumed; $\Phi$\_c join-dominant \\
tensor & \texttt{Hv1\_open \otimes 2GBI} & T+P CONFLICT; inhibitor occludes, not merges \\
path & \texttt{AtHv1\_silent -\textgreater{} AtHv1\_primed} & BLOCKED; priming is a discrete T-class jump \\
cofactor & \texttt{cofactor(cooper\_pair, electron)} & Reconstructs phonon-dressed partner \\
pipeline & \texttt{Hv1\_open \textgreater{}\textgreater{} meet \textgreater{}\textgreater{} join} & $\Delta$\xi=0; $\Phi$\_c throughout; monad-law compliant \\
\bottomrule
\end{tabularx}
\end{table}

The d(AtHv1\_primed, PsHv1) = 0.000 result --- the gymnosperm channel is algebraically identical to mechanically primed Arabidopsis --- is the single number that collapses the phylogenetic argument. The RSN (Ring-Shaped Network, K\_trap) is the \emph{entire} primitive distinction between angiosperm and gymnosperm Hv channels.

\begin{center}
\rule{0.5\textwidth}{0.4pt}
\end{center}

\begin{enumerate}
    \item \textbf{\texttt{mplus} semantics on cost}: should the fallback branch's cost be added to the primary branch's failed cost, or reset? Current design: costs accumulate (Writer is append-only). Alternative: reset on fallback (cleaner but loses failed-branch evidence).
    \item \textbf{Context merging on \texttt{bind}}: \texttt{F-floor} takes max (strictest requirement wins). Is this always right for \texttt{meet} operations which lower F? Currently: \texttt{meet} does not update \texttt{context.f\_floor} (it lowers F, not raises the floor). Only \texttt{join} and \texttt{HotSwap} raise the floor.
    \item \textbf{\texttt{assert} timeout}: criticality probe assertions (\texttt{phi\_c\_score \textgreater{} N}) run the full Varma probe. For large catalogs this could be slow in a pipeline. Add \texttt{--timeout} flag or cache probe results per synthon.
    \item \textbf{Named strategies in \texttt{.syn}}: the \texttt{strategies:} block in \texttt{.syn} allows defining reusable sub-pipelines. Should these be first-class \texttt{DesignStrategy} objects (composable with \texttt{\textgreater{}\textgreater{}} and \texttt{\textless{}\textbar{}\textgreater{}}) or just YAML macros (inlined at parse time)?
    \item \textbf{Phase boundary semantics}: the Ward dendrogram gives an empirical boundary at d $\approx$ 9.52 for the current 8-synthon quantum catalog. Does this threshold generalise? Adding classical synthons will compress the inter-cluster gap. The open question is whether a \textbf{universal} threshold exists in the Kleisli metric, or whether boundaries are always catalog-relative. Phase 3e should test this by mixing the quantum catalog with the full classical molecular catalog and measuring whether the topological/particle split persists at the same threshold.
    \item \textbf{\texttt{.syn} assert grammar extension}: predicates for $\Omega$ and quantum grammar are not yet in the dispatch table. Needed additions: \texttt{topo\_index == Omega\_NA}, \texttt{grammar == QUANTUM\_AND}, \texttt{K\_MBL}, \texttt{factor8\_triggered}. These enable design programs that assert topological protection before quantum tensor operations.
\end{enumerate}


\end{document}
